% ------------------------------------------------------------
% Page 80
% ------------------------------------------------------------

\begin{center}
    {\Large\bfseries
    LA SOCIETE CEVENOLE DES CARBURANTS\\
    FORESTIERS
    }

    \vspace{0.7cm}

    {\Large\bfseries
    A MEYRUEIS
    }

    \vspace{0.9cm}

    Témoignage sur son action pendant l’occupation\\
    par Robert Cabanel, ex – comptable à l’Usine
\end{center}

\vspace{1.1cm}

La pénurie d’essence, à cette époque, avait orienté les recherches en carburants de
remplacement. Le charbon de bois faisait tourner les moteurs à gazogène. Un brevet suisse,
exploité en usine à Chapareillan près de Grenoble ayant donné de bons résultats, incita les
dirigeants à monter d’autres usines. C’est ainsi qu’après avoir visité le massif forestier de
l’Aigoual, les dirigeants de l’usine de Chapareillan se mirent en rapport avec un groupe
financier de Nîmes ayant à sa tête la banque Arnaud et Gaidan et décidèrent de monter une
usine à Meyrueis. La Société Cévenole des Carburants Forestiers fut constituée avec Meyrueis
comme siège social.

\vspace{0.45cm}

\textbf{Les débuts de l’Usine d’Alauze}

\vspace{0.35cm}

La construction démarra en hiver 1942 sur un terrain situé au lieu dit Alauze, à deux km de
Meyrueis dans la vallée de la Brèze. La forêt était à une quinzaine de km. Il s’agissait donc de
fabriquer du charbon de bois en usine dans des fours : les sous – produits, goudrons et jus
pyrolygneux étaient récupérés et, après une certaine préparation, on arrivait à fabriquer des
agglomérés de charbon de bois sous forme d’ovoïdes dont le rendement dans les gazos était
supérieur au charbon de bois ordinaire.

\vspace{0.45cm}

Je dois dire que j’étais, depuis le 6 Septembre 1941, Greffier de la Justice de Paix de
Meyrueis.. Situation de rencontre nécessitée par les circonstances du moment. En effet, après
avoir été mobilisé au 104\textsuperscript{o} bataillon de chasseurs alpins j’avais eu la chance, en faisant partie
de l’armée des Alpes, de pouvoir rentrer dans mon foyer à Cannes où j’exerçais depuis 10 ans
la profession d’employé de banque. La situation économique allait en se dégradant. Des bruits
les plus pessimistes circulaient. On parlait d’envoyer bientôt des Français travailler en
Allemagne. Les employés de banque et les employés d’assurance semblaient les plus visés.
Ma femme attendait notre 2\textsuperscript{o} enfant. Je ne voulais à aucun prix partir en Allemagne et
j’acceptais donc la proposition de mon père, retiré à Meyrueis, consistant à acheter le Greffe
de la Justice de Paix qui allait être vacant et à venir travailler quelques parcelles de notre
petite propriété de Pradines. Après les démarches administratives assez longues, j’obtins
l’autorisation de déménager au titre du « Retour à la Terre » préconisé par le régime de
Vichy.

\vspace{0.45cm}

Quitter Cannes pour Meyrueis constituait pour nous un double sacrifice :

En effet j’abandonnais un traitement de 1.800 F par mois pour une indemnité de fonction de
Greffier de 530 F par mois.

D’autre part la côte d’Azur avec la douceur de son climat et la petite villa confortable que
nous habitions n’avaient rien de comparable avec notre maison de campagne de Pradines.

\clearpage
